\chapter{प्रकाश के स्रोत और संचरण}\label{ch:g7:light-sources-propagation}

तुम यह पन्ना पढ़ क्यों पा रहे हो? लैंप चमकता है --- पर पन्ना कोई लैंप
नहीं, फिर भी तुम्हारी आँख को छपी हुई हर पंक्ति से \omterm{def:g1:light-and-shadows:source}{प्रकाश} मिलता है।
\omterm{def:g1:light-and-shadows:source}{प्रकाश} बनाने वाली चीज़ों और उसे बस आगे बढ़ा देने वाली चीज़ों के बीच वह
भेद है जिसके बिना प्रकाशिकी जी ही नहीं सकती; और उसी के साथ इस पुस्तक
के सबसे पुराने सवाल का पूरा उत्तर भी आता है: किसी चीज़ को देखने के लिए
चाहिए क्या?

\section{बनाने वाले और आगे बढ़ाने वाले}

\begin{definition}[प्राथमिक स्रोत]\label{def:g7:light-sources-propagation:primary}
\emph{प्राथमिक स्रोत}\index{प्राथमिक स्रोत} अपना \omterm{def:g1:light-and-shadows:source}{प्रकाश} ख़ुद बनाता है:
सूर्य और तारे, लपटें, दमकते तंतु, भट्ठी का सफ़ेद-गरम इस्पात, जुगनू,
बिजली की कौंध। उसका \omterm{def:g1:light-and-shadows:source}{प्रकाश} तब भी रहता है जब और कुछ न चमके --- दुनिया
की सारी बत्तियाँ बुझा दो, प्राथमिक स्रोत फिर भी दमकता है।
\end{definition}

\begin{definition}[विसरक वस्तुएँ]\label{def:g7:light-sources-propagation:diffusion}
बाक़ी हर दिखने वाली चीज़ \emph{विसरक वस्तु}\index{विसरण} है: वह \omterm{def:g1:light-and-shadows:source}{प्रकाश}
पाती है और उसका कुछ हिस्सा हर दिशा में फिर से भेज देती है --- यानी
\emph{विसरित} कर देती है। पन्ना, सेब, \omterm{def:g2:moon-in-the-sky:moon}{चंद्रमा}, तुम्हारा अपना चेहरा:
इनमें से कोई अपनी एक झलक भी नहीं बनाता; हर एक पाया हुआ \omterm{def:g1:light-and-shadows:source}{प्रकाश} उदारता
से चारों ओर बिखेर देता है। अँधेरे में कोई विसरक बस अनदेखा रह जाता है।
\end{definition}

\begin{omfigure}
\begin{tikzpicture}[scale=0.8]
  % the lamp, a primary source
  \fill[omThm] (0.7,2.7) circle (0.3);
  \foreach \a in {0,45,...,315} {
    \draw[omThm, thick] (0.7,2.7) ++(\a:0.42) -- ++(\a:0.2);
  }
  \node[above, font=\small] at (1.1,3.35) {लैंप: प्राथमिक स्रोत};
  % its light travels to the apple
  \draw[-Stealth, omThm, thick] (1.35,2.75) -- (3.55,2.15);
  \draw[-Stealth, omThm, thick] (1.35,2.45) -- (3.45,1.65);
  \draw[-Stealth, omThm, thick] (1.25,2.15) -- (3.35,1.15);
  % the apple scatters it to every side
  \draw[thick, fill=omMeth, fill opacity=0.5] (5.2,1.2) circle (0.5);
  \foreach \a in {15,60,105,150,195,240,285,330} {
    \draw[-Stealth, omDef, thick] (5.2,1.2) ++({\a}:0.6)
      -- ++({\a}:1.05);
  }
  \node[below, font=\small] at (5.2,-0.5) {सेब: विसरक};
  \node[right, font=\small, omDef] at (7.35,2.15)
    {विसरित प्रकाश हर ओर जाता है:};
  \node[right, font=\small, omDef] at (7.35,1.55)
    {कमरे की हर आँख को कुछ न कुछ मिलता है};
\end{tikzpicture}

{\small काम का बँटवारा: लैंप \omterm{def:g1:light-and-shadows:source}{प्रकाश} बनाता है; और सेब उसे पाकर उसका कुछ
हिस्सा हर दिशा में बिखेर देता है --- इसीलिए मेज़ के चारों ओर बैठा हर
कोई वही सेब देख लेता है।}
\end{omfigure}

\begin{example}[विसरण दर्पण जैसा लौटाना नहीं है]\label{ex:g7:light-sources-propagation:notmirror}
\omterm{def:g5:mirrors-reflection:reflection}{दर्पण} और पन्ना, दोनों पाया हुआ \omterm{def:g1:light-and-shadows:source}{प्रकाश} वापस भेजते हैं --- पर उलटे ढंग
से। \omterm{def:g5:mirrors-reflection:reflection}{दर्पण} उसे \emph{एक} अनुशासित दिशा में लौटाता है (बराबर कोण ---
तुम्हारा पुराना टकराव का नियम) और तसवीर बचा लेता है; जबकि पन्ने का
सूक्ष्म खुरदरापन \omterm{def:g1:light-and-shadows:source}{प्रकाश} को \emph{हर ओर} छिड़क देता है, तसवीर मिटा देता
है, पर ख़ुद को एक साथ हर आँख के सामने पेश कर देता है। इसीलिए \omterm{def:g5:mirrors-reflection:reflection}{दर्पण} में
लैंप दिखता है, जबकि पन्ना तुम जहाँ भी खड़े हो वहाँ से सिर्फ़ ख़ुद को
दिखाता है।
\end{example}

\begin{example}[चंद्रमा की फ़ाइल, बंद]\label{ex:g7:light-sources-propagation:moon}
बचपन वाले फ़ैसले को अब उसके औपचारिक शब्द मिल जाते हैं: \omterm{def:g2:moon-in-the-sky:moon}{चंद्रमा} कोई
\omterm{def:g7:light-sources-propagation:primary}{प्राथमिक स्रोत} नहीं है; वह \omterm{def:g4:solar-system:planet}{ग्रह} के नाप की एक \omterm{def:g7:light-sources-propagation:diffusion}{विसरक वस्तु} है --- सलेटी
चट्टान, जो सूर्य के \omterm{def:g1:light-and-shadows:source}{प्रकाश} को हर दिशा में बिखेरती है, और हमारी ओर भी।
यानी चाँदनी दूसरे हाथ की धूप है, और उसका ``अँधेरा हिस्सा'' बस वह आधा
भाग है जिसकी फ़ाइल में उस पल आगे बढ़ाने के लिए कोई \omterm{def:g1:light-and-shadows:source}{प्रकाश} ही नहीं होता।
\end{example}

\section{देखना, निपट गया}

\begin{proposition}[दृष्टि की दो शर्तें]\label{prop:g7:light-sources-propagation:vision}
आँख किसी वस्तु को ठीक तब देखती है जब
\begin{enumerate}
  \item वस्तु से \omterm{def:g1:light-and-shadows:source}{प्रकाश} \emph{निकले} --- अगर वह \omterm{def:g7:light-sources-propagation:primary}{प्राथमिक स्रोत} है तो
        उसका बनाया हुआ; और अगर वह आगे बढ़ाने वाली है तो विसरित ---
        और
  \item वह \omterm{def:g1:light-and-shadows:source}{प्रकाश} बिना रुकावट, सीधी रेखा में आँख के भीतर पहुँचे।
\end{enumerate}
इनमें से कोई एक शर्त तोड़ दो --- बिना उजाले वाला विसरक, या नज़र की रेखा
पर खड़ी दीवार --- और वस्तु अनदेखी रह जाती है। इस पुस्तक की हर दिखने-न
दिखने वाली पहेली इन्हीं दो जाँचों में खुल जाती है।
\end{proposition}

\begin{example}[दिखने लगा किरणपुंज]\label{ex:g7:light-sources-propagation:beam}
साफ़ \omterm{def:g2:air-around-us:air}{हवा} में कमरा पार करता टॉर्च का \omterm{def:g6:light-propagation:ray}{किरणपुंज} बग़ल से दिखता ही नहीं ---
\omterm{def:g2:air-around-us:air}{हवा} लगभग कुछ भी आगे नहीं बढ़ाती, और पुंज का कोई \omterm{def:g1:light-and-shadows:source}{प्रकाश} तुम्हारी आँख की
ओर मुड़ता नहीं। अब कोई धूल भरा कपड़ा झाड़ो: पुंज उछलकर सामने आ जाता
है। धूल का हर कण एक नन्हा विसरक है, जो पुंज का \omterm{def:g1:light-and-shadows:source}{प्रकाश} पकड़कर उसका कुछ
हिस्सा तुम्हारी ओर बिखेर देता है --- हज़ारों कण मिलकर तुम्हारे लिए पुंज
की सीधी रेखा खींच देते हैं। कोहरा, धुंध और धुआँ लाइटहाउस के पुंजों और
सिनेमा के प्रक्षेपकों के लिए यही सेवा करते हैं।
\end{example}

\section{प्रकाश के लिए तीन लिबास}

\begin{definition}[पारदर्शी, पारभासी, अपारदर्शी]\label{def:g7:light-sources-propagation:coats}
पदार्थ \omterm{def:g1:light-and-shadows:source}{प्रकाश} से तीन ढंग से मिलते हैं:
\begin{enumerate}
  \item \emph{पारदर्शी}\index{पारदर्शी}: \omterm{def:g1:light-and-shadows:source}{प्रकाश} लगभग बिना छेड़े,
        अपना सलीक़ा बनाए हुए पार निकल जाता है --- साफ़ काँच, ठहरा
        पानी, \omterm{def:g2:air-around-us:air}{हवा}; \omterm{prop:g5:mirrors-reflection:image}{प्रतिबिंब} सही-सलामत पार जाते हैं;
  \item \emph{पारभासी}\index{पारभासी}: \omterm{def:g1:light-and-shadows:source}{प्रकाश} पार तो जाता है, पर
        गड्डमड्ड होकर --- धुँधला काँच, ट्रेसिंग काग़ज़, पतला परदा;
        चमक पार जाती है, \omterm{prop:g5:mirrors-reflection:image}{प्रतिबिंब} नहीं;
  \item \emph{अपारदर्शी}\index{अपारदर्शी}: \omterm{def:g1:light-and-shadows:source}{प्रकाश} रोक लिया जाता है
        --- या तो सोख लिया जाता है या वापस विसरित कर दिया जाता है
        --- लकड़ी, पत्थर, धातु, तुम्हारा हाथ; और अपारदर्शी चीज़ के
        पीछे \omterm{def:g1:light-and-shadows:shadow}{छाया} पड़ती है।
\end{enumerate}
\end{definition}

\begin{method}[टॉर्च की रोशनी में पदार्थ छाँटना]\label{met:g7:light-sources-propagation:sort}
एक अँधेरा कमरा, एक टॉर्च, एक छपा हुआ पन्ना, और बीच में नमूना:
\begin{enumerate}
  \item नमूने को छपाई के ऊपर पकड़ो और टॉर्च उसके पीछे रखो;
  \item उसमें से छपाई पढ़ी जा रही है? --- \omterm{def:g7:light-sources-propagation:coats}{पारदर्शी};
  \item दमक दिखती है पर अक्षर नहीं? --- \omterm{def:g7:light-sources-propagation:coats}{पारभासी};
  \item अँधेरा दिखता है, और उस पार \omterm{def:g1:light-and-shadows:shadow}{छाया}? --- \omterm{def:g7:light-sources-propagation:coats}{अपारदर्शी}।
\end{enumerate}
ईमानदारी से जाँचो: बहुत-से पदार्थ मोटाई के साथ पाला बदल लेते हैं ---
पानी की एक परत \omterm{def:g7:light-sources-propagation:coats}{पारदर्शी} है, पर गहरा समुद्र रात जैसा काला; काग़ज़ की एक
शीट \omterm{def:g7:light-sources-propagation:coats}{पारभासी} है, पर बंद किताब \omterm{def:g7:light-sources-propagation:coats}{अपारदर्शी}।
\end{method}

\begin{omfigure}
\begin{tikzpicture}[scale=0.8]
  \foreach \x/\name in {0.5/पारदर्शी, 4.4/पारभासी,
      8.3/अपारदर्शी} {
    \fill[omThm] (\x,1.5) circle (0.22);
    \draw[thick] (\x+1.3,0.4) rectangle (\x+1.7,2.6);
    \node[below, font=\small] at (\x+1.6,-0.0) {\name};
  }
  \draw[-Stealth, omThm, thick] (0.75,1.5) -- (1.78,1.5);
  \draw[-Stealth, omThm, thick] (2.22,1.5) -- (3.4,1.5);
  \draw[-Stealth, omThm, thick] (4.65,1.5) -- (5.68,1.5);
  \foreach \a in {28,10,-12,-30} {
    \draw[-Stealth, omThm, thin] (6.12,1.5) -- ++({\a}:1.1);
  }
  \draw[-Stealth, omThm, thick] (8.55,1.5) -- (9.58,1.5);
  \draw[line width=3pt, omExo] (10.02,0.55) -- (10.02,2.45);
  \node[right, font=\footnotesize, omExo] at (10.1,0.8) {छाया वाली ओर};
\end{tikzpicture}

{\small किसी किरण के तीन अंजाम: सलीक़े के साथ पार; पार, पर गड्डमड्ड
होकर; और रुक जाना, पीछे \omterm{def:g1:light-and-shadows:shadow}{छाया} के साथ।}
\end{omfigure}

\begin{example}[वास्तुकला की तीनों टोलियाँ]\label{ex:g7:light-sources-propagation:architecture}
राजगीर तीनों को जान-बूझकर तैनात करते हैं। \omterm{def:g7:light-sources-propagation:coats}{पारदर्शी} खिड़की का काँच:
नज़ारा और \omterm{def:g1:light-and-shadows:source}{प्रकाश}, दोनों। \omterm{def:g7:light-sources-propagation:coats}{पारभासी} स्नानघर के शीशे और धुँधले दफ़्तरी
पर्दे: दिन का उजाला भीतर, और निजता बनी हुई --- \omterm{prop:g5:mirrors-reflection:image}{प्रतिबिंब} के बिना चमक
देना ही उनका काम है। \omterm{def:g7:light-sources-propagation:coats}{अपारदर्शी} दीवारें और किवाड़: माँगने पर अँधेरा।
लैंप के आवरण आराम की सेवा में लगी पारभासिता हैं: दहकते तंतु की चौंध
गड्डमड्ड होकर मुलायम, बिना स्रोत वाली दमक बन जाती है।
\end{example}

\begin{remark}[आगे बढ़ाने को कुछ नहीं, तो देखने को कुछ नहीं]\label{rem:g7:light-sources-propagation:space}
पहेली: तारों के बीच \omterm{def:g1:floating-and-sinking:words}{तैरती} एक अंतरिक्ष-यात्री धूप से भरे शून्य के पार
\emph{बग़ल} में देखती है। उसे क्या दिखता है? कालापन --- जबकि सूर्य का
\omterm{def:g1:light-and-shadows:source}{प्रकाश} उसकी आँखों के पास से बहता जा रहा है। अंतरिक्ष में गिनने लायक
धूल है ही नहीं: गुज़रते \omterm{def:g1:light-and-shadows:source}{प्रकाश} को उसकी ओर विसरित करने वाला कुछ नहीं,
और जो \omterm{def:g1:light-and-shadows:source}{प्रकाश} आँख में न घुसे वह कुछ भी नहीं रँगता। दिन का आकाश नीला
ठीक इसीलिए है कि हमारी \omterm{def:g2:air-around-us:air}{हवा} गुज़रती धूप को हमारी ओर आगे \emph{बढ़ाती}
है --- यानी सिर के ऊपर एक पूरे \omterm{def:g4:solar-system:planet}{ग्रह} भर का कोमल \omterm{def:g7:light-sources-propagation:diffusion}{विसरण}। \omterm{def:g2:air-around-us:air}{हवा} नहीं, तो
आकाश काला, दोपहर में भी: \omterm{def:g2:moon-in-the-sky:moon}{चंद्रमा} के अंतरिक्ष-यात्री धूप में खड़े थे और
ऊपर तारों-रहित काला गुंबद था।
\end{remark}

\section{अभ्यास}

\begin{exercise}[$\star$]\label{exo:g7:light-sources-propagation:1}
छाँटो: मोमबत्ती की लौ; \ \omterm{def:g2:moon-in-the-sky:moon}{चंद्रमा}; \ जुगनू; \ यह पन्ना; \ यातायात की
लाल बत्ती; \ साँझ का शुक्र \omterm{def:g4:solar-system:planet}{ग्रह}।
\end{exercise}

\begin{exercise}[$\star$]\label{exo:g7:light-sources-propagation:2}
कोई \omterm{def:g7:light-sources-propagation:diffusion}{विसरक वस्तु} मिले हुए \omterm{def:g1:light-and-shadows:source}{प्रकाश} के साथ क्या करती है --- और अँधेरे कमरे
में उसका क्या होता है?
\end{exercise}

\begin{exercise}[$\star$]\label{exo:g7:light-sources-propagation:3}
दृष्टि की दोनों शर्तें बताओ, और उन्हीं से समझाओ: बिलकुल साफ़ \omterm{def:g5:mirrors-reflection:reflection}{दर्पण} की
\emph{अपनी} सतह क्यों नहीं दिखती, सिर्फ़ उसमें दिखने वाली चीज़ें ही
क्यों दिखती हैं? (ध्यान से --- बारीक सवाल है; उत्तर सीधा दो: उसकी
चमकाई हुई सतह \omterm{def:g1:light-and-shadows:source}{प्रकाश} के साथ क्या करने से इनकार कर देती है?)
\end{exercise}

\begin{exercise}[$\star$]\label{exo:g7:light-sources-propagation:4}
मेज़ के चारों ओर बैठा हर कोई वही सेब क्यों देख लेता है, जबकि चम्मच में
लैंप का \omterm{prop:g5:mirrors-reflection:image}{प्रतिबिंब} सिर्फ़ एक ही ख़ुशक़िस्मत जगह पर बैठी आँख को दिखता है?
\end{exercise}

\begin{exercise}[$\star$]\label{exo:g7:light-sources-propagation:5}
साफ़ \omterm{def:g2:air-around-us:air}{हवा} में टॉर्च का पुंज बग़ल से क्यों नहीं दिखता, और धुंध में क्यों
दिख जाता है? इस प्रक्रिया का नाम बताओ।
\end{exercise}

\begin{exercise}[$\star$]\label{exo:g7:light-sources-propagation:6}
वर्गीकृत करो: खिड़की का काँच; \ ट्रेसिंग काग़ज़; \ एक ईंट; \ उथला
साफ़ पानी; \ धुँधले \omterm{def:g3:first-electric-circuit:bulb}{बल्ब} का ख़ोल; \ एल्यूमिनियम की पन्नी।
\end{exercise}

\begin{exercise}[$\star$]\label{exo:g7:light-sources-propagation:7}
स्नानघर की खिड़की और लैंप के आवरण में पारभासिता कौन-सा काम करती है?
क्या पार जाता है, और क्या रोक लिया जाता है?
\end{exercise}

\begin{exercise}[$\star$]\label{exo:g7:light-sources-propagation:8}
\omterm{def:g2:moon-in-the-sky:moon}{चंद्रमा} पर धूप पड़ती है --- फिर भी वहाँ गए अंतरिक्ष-यात्रियों को
दोपहर में काला आकाश दिखा।
\cref{rem:g7:light-sources-propagation:space} से समझाओ, और बताओ कि
पृथ्वी की दोपहर का आकाश इसके बजाय चमकीला नीला क्यों होता है।
\end{exercise}

\begin{exercise}[$\star\star$]\label{exo:g7:light-sources-propagation:9}
टॉर्च के सामने काग़ज़ की एक शीट गरम-सी दमक उठती है; पर बंद किताब
\omterm{def:g1:light-and-shadows:source}{प्रकाश} पूरी तरह रोक देती है। मोटाई किसी पदार्थ की टोली-सदस्यता के साथ
क्या करती है? पाला बदलने वाला एक और उदाहरण दो।
\end{exercise}

\begin{exercise}[$\star\star$]\label{exo:g7:light-sources-propagation:10}
सिनेमा की तरकीब: साफ़ हॉल में प्रक्षेपकों के पुंज दिखते ही नहीं, फिर
भी फ़िल्मों में वे \omterm{def:g2:air-around-us:air}{हवा} को नाटकीय ढंग से चीरते दिखाए जाते हैं। सेट की
टोली को \omterm{def:g2:air-around-us:air}{हवा} में क्या जोड़ना पड़ता है, और वे कौन-सी भौतिकी ख़रीद रहे
हैं?
\end{exercise}

\begin{exercise}[$\star\star$]\label{exo:g7:light-sources-propagation:11}
अलाव के उजाले में तुम्हें किसी साथी का चेहरा दिखता है। \omterm{def:g1:light-and-shadows:source}{प्रकाश} की पूरी
यात्रा बताओ --- स्रोत, \omterm{def:g7:light-sources-propagation:diffusion}{विसरण}, सीधी रेखाएँ --- और गिनो कि अगर तुम्हें
वही चेहरा झील में भी दिखे, तो शृंखला में कितने विसरक खड़े हैं।
\end{exercise}

\begin{exercise}[$\star\star\star$]\label{exo:g7:light-sources-propagation:12}
प्रकृति का अपना प्रकाश-तमाशा: चमकीले दिन में सूर्य से दूर का आकाश नीला
होता है और सूर्य की थाली चौंधियाती सफ़ेद --- और \omterm{def:g1:day-and-night:daynight}{सूर्यास्त} पर नीचा
सूर्य नारंगी हो जाता है जबकि ऊपर के बादल गुलाबी हो उठते हैं। सिर्फ़
इतना मानकर कि ``\omterm{def:g2:air-around-us:air}{हवा} अपने में से गुज़रते \omterm{def:g1:light-and-shadows:source}{प्रकाश} का थोड़ा हिस्सा विसरित
कर देती है, और रास्ता जितना लंबा उतना ज़्यादा'', समझाओ कि नीचे उतरा
सूर्य अपने \omterm{def:g1:light-and-shadows:source}{प्रकाश} का एक हिस्सा क्यों गँवा देता है --- और चुराया गया वह
हिस्सा किसे मिलता है। (रंगों की पूरी कहानी अगले साल के वर्णक्रम वाले
अध्याय में है; यहाँ सिर्फ़ ज़्यादा-रास्ता, ज़्यादा-विसरण से तर्क करो।)
\end{exercise}

\section{समस्या: संग्रहालय को रोशन करना}

\begin{problem}[{सप्ताहांत समस्या --- नया मूर्ति-कक्ष; प्रकाश-डिज़ाइनर
की रात की पाली; स्रोत, विसरक और तीन तरह के काँच}]\label{pb:g7:light-sources-propagation:1}
संग्रहालय का नया कक्ष सोमवार को खुलता है: संगमरमर की एक मूर्ति,
दीवारों पर चित्र, काँच की एक प्रदर्शन-पेटी, और एक घबराया हुआ
संग्रहाध्यक्ष। तुम प्रकाश-डिज़ाइनर की मदद कर रहे हो।

\medskip\noindent\textbf{भाग I --- मौक़े पर पहले सिद्धांत।}
\begin{enumerate}
  \item संग्रहाध्यक्ष पूछते हैं कि ``इतनी चमकदार सफ़ेद'' मूर्ति को
        लैंपों की ज़रूरत ही क्यों है। दृष्टि की दोनों शर्तों से उत्तर
        दो।
  \item लैंप जलने के बाद कक्ष के \omterm{def:g7:light-sources-propagation:primary}{प्राथमिक स्रोत} और विसरक गिनाओ: लैंप,
        मूर्ति, चित्र, दर्शक, और पेटी का काँच।
  \item संगमरमर हर ओर से दमकता-सा लगता है, और कहीं कोई चौंधियाता
        धब्बा नहीं। \omterm{def:g1:light-and-shadows:source}{प्रकाश} लौटाने का कौन-सा ढंग यहाँ काम कर रहा है,
        और वह दर्शकों की हर जगह तक क्यों पहुँच जाता है?
  \item इसके उलट, चमकाई हुई पीतल की एक तख़्ती किसी लैंप का एक
        चौंधियाता \omterm{prop:g5:mirrors-reflection:image}{प्रतिबिंब} फेंकती है। दूसरे ढंग का नाम बताओ, और वह
        नियम भी जो तय करता है कि वह चौंध जाती कहाँ है।
\end{enumerate}

\medskip\noindent\textbf{भाग II --- चौंध को साधना।}
\begin{enumerate}[resume]
  \item नंगे \omterm{def:g3:first-electric-circuit:bulb}{बल्ब} दर्शकों की आँखों में चुभते हैं। डिज़ाइनर हर एक को
        धुँधले गोले में ढक देती है। यह गोला किस टोली का है, और वह
        तंतु के \omterm{def:g1:light-and-shadows:source}{प्रकाश} के साथ क्या करता है?
  \item चित्रों पर \omterm{def:g1:light-and-shadows:source}{प्रकाश} पड़ना चाहिए, पर मुख्य दर्शक-स्थान से उनकी
        वार्निश में कोई लैंप \emph{प्रतिबिंबित} नहीं दिखना चाहिए।
        \omterm{def:g5:mirrors-reflection:reflection}{दर्पण} वाले नियम से बताओ कि डिज़ाइनर लैंप चित्र और दर्शक के
        मुक़ाबले कहाँ रखती है।
  \item प्रदर्शन-पेटी का साफ़ काँच लगभग अनदेखा रहता है --- जब तक उसमें
        किसी दर्शक की सफ़ेद आस्तीन किसी साये की तरह न उभर आए। साफ़
        काँच पार भी क्यों जाने देता है और हल्का-सा \omterm{def:g5:mirrors-reflection:reflection}{दर्पण} भी क्यों बन
        जाता है? पेटी के पीछे डिज़ाइनर की बनाई काली दीवार वाली ताक़
        क्या हासिल करती है?
  \item कभी-कभी धूल के कण स्पॉटलाइट के पुंजों में से बहते हैं और \omterm{def:g2:air-around-us:air}{हवा}
        में चमकीले धागे खींच देते हैं। संग्रहाध्यक्ष: ``कितना सुंदर!
        इसे रहने दो।'' डिज़ाइनर: ``इसका मतलब है गंदी \omterm{def:g2:air-around-us:air}{हवा}।'' भौतिकी के
        हिसाब से सही कौन है, और वे धागे हैं क्या?
\end{enumerate}

\medskip\noindent\textbf{भाग III --- रात की जाँच।}
\begin{enumerate}[resume]
  \item बत्तियाँ बंद, किवाड़ सील: पहरेदार ख़बर देता है कि कक्ष
        ``बिलकुल काला है --- वह सफ़ेद मूर्ति तक नहीं दिखती''। दोनों
        शर्तों से समझाओ कि अँधेरे में सफ़ेदी क्यों बेबस है।
  \item ऊपर धुँधले काँच का एक रोशनदान प्रस्तावित है। वह दिन में क्या
        देगा, और क्या देने से इनकार कर देगा (बादलों का नज़ारा, सीधी
        धूप की चौंध)?
  \item संग्रहाध्यक्ष को डर है कि धूप चित्रों के रंग उड़ा देगी, और वे
        ``बिना किसी सूर्य वाला \omterm{def:g1:light-and-shadows:source}{प्रकाश}'' माँगते हैं। डिज़ाइनर मुस्कुरा
        देती है: दिन में तो उत्तर की खिड़कियों का मुलायम \omterm{def:g1:light-and-shadows:source}{प्रकाश} भी
        धूप \emph{ही} है। समझाओ --- उसे आगे बढ़ाया किसने?
  \item डिज़ाइनर की रिपोर्ट लिखो: तीन वाक्य --- एक इस पर कि \omterm{def:g7:light-sources-propagation:diffusion}{विसरण}
        संग्रहालय का सबसे अच्छा दोस्त क्यों है, एक इस पर कि \omterm{def:g5:mirrors-reflection:reflection}{दर्पण}
        वाले परावर्तनों पर कहाँ नज़र रखनी चाहिए, और एक इस पर कि हर
        ज़िम्मेदारी के लिए काँच की कौन-सी टोली चुनी गई।
\end{enumerate}
\end{problem}
